% !TeX root = ../main.tex
\section{Audio primer}
\label{sec:audio}

This section introduces the signal-processing ideas used everywhere else. It
assumes nothing beyond high-school mathematics.

\subsection{Samples, frames and the spectrogram}

A digital recording is a list of numbers, the \term{samples}, measured at a
fixed \term{sample rate}. The app's audio engine runs at 48\,000 samples per
second (48\kHz); the models resample to what they were trained on (16\kHz for
the guitar CRNN, 22.05\kHz for Basic Pitch and Beat This!, 44.1\kHz stereo for
BS-RoFormer).

Pitch is easier to see in frequency than in time. The \term{short-time Fourier
transform} (STFT) cuts the signal into overlapping \term{frames} of $N$ samples
(the \term{window}, e.g.\ 2048), moves forward by a \term{hop} of $H$ samples
(e.g.\ 512) between frames, and computes the spectrum of each frame: how much
energy there is at each frequency. Stacking these spectra side by side gives a
\term{spectrogram}, an image with time on the horizontal axis, frequency on the
vertical axis and energy as brightness. Most transcription models take a
spectrogram as input.

\begin{worked}[: resolution]
At 16\kHz, a hop of 160 samples gives $16000/160 = 100$ frames per second, so
time is resolved in steps of 10\ms. A window of $N$ samples resolves frequency
in steps of $f_s/N$: with $N=2048$ at 16\kHz that is 7.8\Hz, less than a
semitone even at the lowest guitar note (E2 = 82.4\Hz; the next semitone, F2, is
87.3\Hz). Longer windows give finer frequency but blur time. This trade-off is
why some systems compute two resolutions.
\end{worked}

Three frequency axes are common:
\begin{itemize}
\item \textbf{Linear} (the plain STFT): equal spacing in Hz.
\item \textbf{Mel}: spacing that follows human pitch perception, roughly linear
  below 1\kHz and logarithmic above. Beat This! uses 128 mel bands.
\item \textbf{Constant-Q transform} (CQT): bins spaced by a fixed musical
  interval (for example a third of a semitone), so every octave has the same
  number of bins. Natural for music; TabCNN and Basic Pitch use CQT-like inputs.
\end{itemize}
Energies are almost always compressed with a logarithm (decibels), because
loudness perception is roughly logarithmic and it keeps quiet notes visible.

\subsection{Pitch, MIDI numbers and harmonics}
\label{sec:audio-harmonics}

Musical pitch is labelled here with \term{MIDI numbers}: A4 = 440\Hz{} is 69,
and each semitone adds one. The frequency of MIDI pitch $p$ is
\[
  f(p) = 440 \cdot 2^{(p-69)/12}\ \text{Hz}.
\]
Standard guitar tuning is 40 (E2), 45, 50, 55, 59, 64 (E4) from string 6 to
string 1; each fret adds one semitone.

A plucked string does not vibrate at one frequency. It produces a
\term{fundamental} $f_0$ plus \term{harmonics} at $2f_0, 3f_0, 4f_0,\ldots$ The
2nd harmonic is one octave up, the 3rd an octave and a fifth, the 4th two
octaves. Every note detector must decide whether energy at $2f_0$ is a harmonic
of a note at $f_0$ or a separate note an octave higher. Getting this wrong
produces an \term{octave error}.

\begin{keyidea}[: why distortion causes octave-low errors]
A distortion circuit clips the waveform. Clipping a sum of two notes at
frequencies $f_1$ and $f_2$ creates energy at $f_2 - f_1$, $2f_1 - f_2$ and
similar combinations (\term{intermodulation}). For a power chord (a root and the
fifth above it, ratio 3:2) the difference tone $f_2 - f_1 = f_1/2$ is one
octave \emph{below} the root. A model trained on clean guitar sees a strong
partial an octave below and reports a note there. This is the main failure we
measure on the ``Metal High Gain'' preset (Section~\ref{sec:bench-app}): most
octave errors are an octave low.
\end{keyidea}

\subsection{Onsets and the pitch contour}

An \term{onset} is the moment a note starts: a burst of broadband energy from
the pick, followed by the harmonic tone. Onset detectors look for sudden rises
in energy across frequencies (\term{spectral flux}: the sum of positive
frame-to-frame increases of the spectrum). A \term{pitch contour} is the
fundamental frequency over time at fine resolution (a fraction of a semitone).
A bend appears in it as a smooth rise of one or two semitones, vibrato as a
periodic wobble of a quarter to half a semitone at 4--8\Hz, a slide as a
staircase of semitone steps. Section~\ref{sec:techniques} builds technique
detection on this.

\subsection{Stereo and panning}

A stereo recording has a left and a right channel. Mixing engineers
\term{pan} each instrument: a guitar panned hard left is only in the left
channel, one in the centre is equal in both. Rock and metal records often place
two rhythm guitars hard left and hard right and a lead guitar in the centre. The
ratio of left to right energy in a time-frequency bin therefore tells roughly
where in the stereo image the sound came from. The app uses this to split
guitars that a separator returns as one stem (Section~\ref{sec:pan-split}).
